\begin{tikzpicture}[scale=1]
\tikzmath{
\L=2.5;
\v=1.5;
}
\scalebox{1.0}{

 \coordinate (P1) at (1.1*\L,0.7*\L);
 \coordinate (P2) at ($(P1)+(-0.5*\L,-0.5*\L)$);
 
 %axes:
 \draw[-latex, line width=1.5pt] (0,0) -- (1.2*\L,0) node[below]{$x$}; 
 \draw[-latex, line width=1.5pt] (0,0) -- (0,\L) node[left]{$y$}; 
 
 \draw[-latex,line width=2pt] (P1) -- ($(P1)+(150:\v)$) node[midway,below]{$\vec v$};
 \draw[dashed,line width=1pt,color=red] ($(P1)+(150:\v)$) -- ($(P1)+(150:\v)+cos(30)*(\v,0)$) node[midway, above] {$v_x$};
 \draw[dashed,line width=1pt,color=red] ($(P1)+(150:\v)+cos(30)*(\v,0)$) --(P1) node[midway, right] {$v_y$};
 
 \draw[-latex,line width=2pt] (P2) -- ($(P2)+(150:\v)$) node[midway,below]{$\vec v$};
 \draw[dashed,line width=1pt,color=red] (P2) -- ($(P2)+(0:\v)$);
 \draw[dashed,line width=1pt,color=red] ($(P2)+(0.2*\v,0)$) arc (0:150:0.2*\v) node[midway, above] {$\phi$};;
}
\end{tikzpicture}
\captionof*{figure}{\textit{A vector in 2D can be specified by two Cartesian components ($v_x,v_y$) or its length, $v$, and an angle with respect to the $x$ direction, $\phi$.}}